\section{What the special cases teach a general local solver}
\label{sec:signed-star-general-lessons}

The special cases support a design principle more precise than ``momentum is
good'' or ``momentum is bad.''  A proved accelerated local solver needs four
separate ledgers.

\begin{enumerate}[leftmargin=*]
\item \textbf{Spectral progress.}  The numerical state must contract the
      $\alpha$-slow modes in $O(1/\sqrt\alpha)$ stages.
\item \textbf{Trajectory locality.}  The total degree volume touched by all
      intermediate states must be bounded, not only the final support.
\item \textbf{State continuation.}  When the active face changes, useful
      accelerated state must be transported or its reset must be amortized.
\item \textbf{Semantic certification.}  The stopping interface must imply the
      requested PPR error without imposing a substantially stronger condition
      that erases the speedup.
\end{enumerate}

Optimal SOR closes all four on the promised star: the parameter is spectrally
optimal, the support is the output-scale star, there are no face changes, and
\Cref{lem:signed-star-semantic-identity} replaces a conservative certificate.
Each negative example deletes one of these properties while retaining some of
the others.

\subsection{Why the exact cancellation is special}

On the star, the normalized adjacency has only two seed-visible eigenmodes.
A center update moves them by equal magnitudes, and the PPR error combines the
two residual layers through the exact Green-function weights
$(1,c_\alpha)$.  At optimal relaxation these weights produce
\begin{equation}
  1-c_\alpha\lambda
  =\frac{2\sqrt\alpha}{1+\alpha},
  \label{eq:signed-star-special-cancellation}
\end{equation}
which cancels the critically damped block index.

On a general fixed support, the semantic error at coordinate $u$ is the dual
quantity
\begin{equation}
  \frac{(\pi-p)_u}{d_u}
  =\left\langle
      r,\,H^{-T}\frac{e_u}{d_u}
    \right\rangle.
  \label{eq:signed-star-dual-semantic-error}
\end{equation}
The relevant Green vector depends on $u$ and need not align with a
two-dimensional invariant subspace.  Absolute-value residual certification
replaces this signed inner product by a worst-case positive envelope and can
therefore lose cancellation.  Computing every exact Green vector, however,
is itself an inverse primitive and cannot be treated as a free local test.

This identifies a concrete generalization target:
\begin{quote}
Construct a locally maintainable upper bound on the dual quantities
\eqref{eq:signed-star-dual-semantic-error} that preserves enough signed
low/high-mode cancellation to stop after $\widetilde O(1/\sqrt\alpha)$ local
stages, while charging every response update and query.
\end{quote}

\subsection{A plausible fixed-face proof route}

For a fixed principal subproblem on an exposed face $S$, let
\(\overline H_S=D_S^{-1/2}H_{SS}D_S^{1/2}\), let $e_S$ be its
principal-subproblem error, and split the spectral space at a threshold
$\tau$ into low and high bands.  The fixed-face energy identity already gives
a pathwise high-mode budget:
\begin{equation}
  \|P_{\ge\tau}D_S^{-1/2}e_S\|_2^2
  \le\frac{2\Phi_S(e_S)}{\tau}.
  \label{eq:signed-star-high-mode-budget}
\end{equation}
The star proof is the two-eigenvalue endpoint of this inequality.  A general
fixed-face theorem would need two additional lemmas:

\begin{enumerate}[leftmargin=*]
\item a local or structurally bounded representation of the low-band dual
      response in \eqref{eq:signed-star-dual-semantic-error};
\item a summable charge showing that any safeguard correction caused by the
      high band is paid by the decrease of
      \eqref{eq:signed-star-high-mode-budget}.
\end{enumerate}

This is more promising than asking for momentum never to trigger a safeguard.
FISTA's leaf-star demonstrates why a universal no-overshoot statement is
implausible near zero complementarity margin.  The goal should instead be a
bounded transient-debt theorem.

\subsection{A plausible changing-face algorithmic route}

The FISTA and ASPR examples together suggest the following requirements for a
purely local accelerated method:

\begin{enumerate}[leftmargin=*]
\item Use a conservative lower iterate for support discovery, so a transient
      extrapolation cannot expose a high-degree boundary row merely because a
      KKT margin is zero.
\item Retain a signed accelerated state inside the already exposed face rather
      than restarting a high-accuracy inner solve after every single
      admission.
\item When a face grows, zero-pad or transport the state through an exact
      algebraic identity, and charge any retraction to a decrease in a
      declared energy or margin bank.
\item Stop using either a proved dual-response surrogate or a final paid
      conversion to a safe semantic certificate.  Do not apply a strong
      residual test at every intermediate stage unless its extra work is
      included.
\end{enumerate}

This is a design program, not a completed theorem.  In particular, a
convergence rate for the accelerated state does not bound the volume of the
conservative discovery state, and a safe discovery state does not by itself
amortize the transported momentum.

\subsection{Where the project-wide proof still stops}

The current graph-uniform AESP--LOCSOR promotion gate requires a
graph-independent early-locality or alternative burn-in bound.  The star
result does not control that quantity: its entire graph already has
$\Theta(1/\varepsilon_{\rm ppr})$ volume, and its face never changes.  Nor
does it settle long-spider behavior, hub-rich transient exploration, or the
cost of maintaining a general response.

The correct promotion boundary is therefore:
\begin{equation}
  \boxed{
  \begin{array}{l}
  \text{Proved: accelerated semantic SOR and matching signed-class scale on
  the center-star.}\\[2pt]
  \text{Open: a fully charged graph-uniform algorithm with the same product
  scale.}
  \end{array}}
  \label{eq:signed-star-promotion-boundary}
\end{equation}

\subsection{Further special families worth testing}

The exact proof should extend whenever the seed-visible dynamics close on a
constant-dimensional equitable quotient and block scans cost the quotient
cell volumes.  Natural next cases are complete bipartite graphs with a
cell-uniform seed, double stars with a small hub quotient, and spiders with a
fixed number of identical arms.  For each family the falsifiable test is the
same:
\begin{enumerate}[leftmargin=*]
\item derive the exact quotient recurrence;
\item compute semantic error, rather than only residual magnitude;
\item check whether the polynomial transient carries a compensating factor of
      $\sqrt\alpha$;
\item charge the expansion of the quotient state back to actual vertex-degree
      work and explicit output.
\end{enumerate}
Failure of the third step would identify a genuine family-specific
obstruction; success without the fourth would be only a small-dimensional
algebraic observation, not a local-work theorem.
